\documentclass[10pt,journal,compsoc]{IEEEtran}
\usepackage{amsmath,amssymb,amsfonts}
\usepackage{booktabs}
\usepackage{array}
\usepackage{graphicx}
\usepackage[hidelinks]{hyperref}
\newcommand{\MW}{\,\mathrm{MW}}
\newcommand{\MeV}{\,\mathrm{MeV}}
\newcommand{\keV}{\,\mathrm{keV}}
\newcommand{\kA}{\,\mathrm{kA}}
\newcommand{\MA}{\,\mathrm{MA}}
\newcommand{\kV}{\,\mathrm{kV}}
\newcommand{\cm}{\,\mathrm{cm}}
\newcommand{\mm}{\,\mathrm{mm}}
\newcommand{\m}{\,\mathrm{m}}
\newcommand{\T}{\,\mathrm{T}}
\newcommand{\MPa}{\,\mathrm{MPa}}
\newcommand{\MJ}{\,\mathrm{MJ}}
\newcommand{\kJ}{\,\mathrm{kJ}}
\newcommand{\ms}{\,\mathrm{ms}}
\newcommand{\ns}{\,\mathrm{ns}}
\newcommand{\ps}{\,\mathrm{ps}}
\newcommand{\mus}{\,\mu\mathrm{s}}
\newcommand{\GHz}{\,\mathrm{GHz}}
\newcommand{\GW}{\,\mathrm{GW}}
\newcommand{\TW}{\,\mathrm{TW}}

\title{Supplementary Material: Extended Derivations, Engineering Trade
Studies, and Operational Specifications for the SHBT-Graser Aneutronic
p--$^{11}$B Power Plant}
\author{}
\markboth{}{}

\begin{document}
\maketitle

\begin{abstract}
This supplementary document collects the extended theoretical derivations,
engineering trade studies, and complete operating tables that accompany the
main article ``A Commercial SHBT-Graser Aneutronic p--$^{11}$B Fusion Power
Plant'' (\texttt{power.pdf}). It contains: the full macro-burst timing
derivation of the C-band graser driver; the topological operator and
modular-invariance formulation of the holographic boundary channel; the
pellet stoichiometry and injector trade study (light-gas versus
electromagnetic launch); the three-field magnetic stopping-radius and
radiant-flash derivations; the dual-domain tenth-value-layer shielding
calculations; and the six-stage operational sequence with the LANR
aggregation-bus specification. All numerical quantities are consistent with
the 70-gate audit matrix (\texttt{verification\_matrix.json}) of the
companion digital twin \cite{mainpaper,sys1own_twin2026}.
\end{abstract}

\section{Graser Driver Macro-Burst Timing Derivation}
The primary specification requires a time-averaged optical beam power of
$P_{\mathrm{avg}} = 25.0\MW$ at a macro-pulse repetition frequency of
$f_{\mathrm{rep}} = 100.0\,\mathrm{Hz}$. Energy conservation fixes the
optical energy delivered per macro-shot:
\begin{equation}
E_{\mathrm{macro}} = \frac{P_{\mathrm{avg}}}{f_{\mathrm{rep}}}
 = \frac{25.0\times10^{6}}{100.0} = 250.0\kJ.
\label{eq:emacro}
\end{equation}
The driver also specifies a constituent micro-pulse of
$E_{\mathrm{micro}} = 100.0\,\mathrm{J}$ in a coherence window
$\tau_{\mathrm{micro}} = 1.0\ps$, i.e.\ an instantaneous peak of
\begin{equation}
P_{\mathrm{peak}} = \frac{E_{\mathrm{micro}}}{\tau_{\mathrm{micro}}}
 = \frac{100.0\,\mathrm{J}}{1.0\times10^{-12}\,\mathrm{s}} = 100.0\TW.
\label{eq:ppeak}
\end{equation}
Treating \eqref{eq:ppeak} as a $100\,\mathrm{Hz}$ isolated shot would
under-drive the core by a factor of $2{,}500$ ($10\,\mathrm{kW}$ average),
while concentrating \eqref{eq:emacro} into a single $1\ps$ envelope would
demand $2.5\times10^{17}\,\mathrm{W}$, beyond any dielectric breakdown
threshold and deep into relativistic self-focusing
\cite{mourou2018,plasma_mirrors2018}. The contradiction is resolved by an
RF bunch hierarchy on a C-band linac at
$f_{\mathrm{rf}} = 5.712\GHz$:
\begin{equation}
T_{\mathrm{rf}} = \frac{1}{f_{\mathrm{rf}}}
 = 1.750700\times10^{-10}\,\mathrm{s} = 175.070\ps.
\end{equation}
The micro-bunch count per macro-shot is
\begin{equation}
N_{\mathrm{micro}} = \frac{E_{\mathrm{macro}}}{E_{\mathrm{micro}}}
 = \frac{250{,}000}{100} = 2{,}500.
\end{equation}
Under synchronous fundamental-bucket filling ($h=1$) the burst duration is
\begin{equation}
\tau_{\mathrm{burst}} = N_{\mathrm{micro}} T_{\mathrm{rf}}
 = 2500\times 175.070\ps = 437.675\ns,
\end{equation}
(875.350 ns for every-second-bucket filling $h=2$), both comfortably inside
the $1.0\mus$ flat-top of cryogenic C-band klystrons \cite{bes2009}. The
intermediate burst power is
\begin{equation}
P_{\mathrm{burst}} = \frac{E_{\mathrm{macro}}}{\tau_{\mathrm{burst}}}
 = \frac{250.0\times10^{3}}{4.37675\times10^{-7}}
 \approx 571.2\GW,
\end{equation}
with macro duty factor
$D_{\mathrm{macro}} = \tau_{\mathrm{burst}} f_{\mathrm{rep}}
 = 4.37675\times10^{-5}$ and a quiescent cooling interval of
$\approx 9.99956\ms$ between bursts. Peak power is set by the 1 ps
micro-pulse (100 TW), intermediate power by the 437.675 ns burst envelope
(571.2 GW), and average power by the 100 Hz repetition (25.0 MW)
\cite{campri2022,mdpi_fel2022,muggli2021,bnl_atf2020}.

\section{Topological Operators and Modular Invariance}
\subsection{Modular reconstruction cycle}
The complete algebraic transformation across the vacuum transition is
governed by the reconstruction cycle operator
\begin{equation}
\hat{R}_{\mathrm{rerender}} = \hat{T}^{\partial}
  \hat{O}_{\mathrm{excitation}}(\theta)\hat{D}_{\mathrm{derender}}^{\dagger},
\label{eq:rerender}
\end{equation}
where $\hat{D}_{\mathrm{derender}}^{\dagger}$ projects state amplitude onto
the $\eta_D = 23/33$ dark branch,
$\hat{O}_{\mathrm{excitation}}(\theta) = e^{-i\theta\hat{Q}}$ executes the
phase-locked $U(1)$ gauge rotation, and $\hat{T}^{\partial}$ is the unitary
Heegaard--Floer boundary relabeling isometry preserving the boundary $L^2$
norm strictly below the $10^{-122}$ holographic noise floor
\cite{sys1own_shbt2026}. The generator Lie algebras of the active boundary
manifold and the dark sector are denoted $\sigma_k^{(i)}$ and
$\tau_k^{(j)}$, and the macroscopic Stinespring dilation
$V_{\mathrm{unified}}^{\mathrm{macro}}$ carries the joint representation.

The transformation strictly enforces the adiabatic invariance condition
\begin{equation}
\Delta S_A = 0,
\label{eq:adiabatic}
\end{equation}
guaranteeing zero non-unitary quantum-information leakage during state
translation. Topological boundary stability is simultaneously constrained
by the Kojima entropy bounds
\begin{equation}
\mathrm{Ent}(\varphi) \le 10^{20}\cdot\mathrm{Vol}(M),
\qquad
\mathrm{Ent}(\varphi_1) \le [M_1:M]\cdot(\ell_{\mathrm{He}}-1)\cdot\log 3,
\label{eq:kojima}
\end{equation}
where $[M_1:M]$ is the covering index of the boundary manifold and
$\ell_{\mathrm{He}}$ the Heegaard genus of the bounding 3-manifold.

\subsection{Modular closure chain}
Within the canonical WZW affine branch $(k_\ell, k_q, K) = (26,8,312)$,
modular closure invariance is equivalent to the simultaneous vanishing of
the framed topological boundary anomaly tensor $\Delta_{\mathrm{fr}}$, the
induced spacetime metric stress error $E_{\mu\nu}$, and the projected
physical Hamiltonian:
\begin{equation}
[M,S^{\partial}] = [M,T^{\partial}] = 0
\;\Longleftrightarrow\; \Delta_{\mathrm{fr}} = 0
\;\Longleftrightarrow\; E_{\mu\nu} = 0
\;\Longleftrightarrow\; \Pi_{\mathrm{phys}}H_{\mathrm{total}}
   \Pi_{\mathrm{phys}} = 0.
\label{eq:closure}
\end{equation}
If a localized terawatt-scale stress-energy gradient were left
uncompensated, it would rupture the closure chain and drive the boundary
toward topological horizon saturation $\Delta N \ge N_{\mathrm{sat}}$,
destroying holographic coherence and dissipating beam energy back into the
driving electronics. PCSS crowbar anomaly interlocks therefore monitor the
modular residual and fire below the $2.5\ns$ threshold before saturation
can occur \cite{trigatron2018,gaas_pulse2025,neuber2021}.

\subsection{Solid-state suppression}
In solid-state media, propagating megavolt-scale gamma flux induces
Compton scattering and $e^+e^-$ pair production cascading into
$6.92\times10^{8}$ optical-phonon up-conversions per interaction, exceeding
the dissipation capacity of CVD diamond
($\kappa \sim 2200\,\mathrm{W/m\cdot K}$) and quenching NbN/MgB$_2$ HTS
buses ($T_c < 11.79\,\mathrm{K}$). Re-rendering in an ultra-high-vacuum
drift chamber ($P < 10^{-10}\,\mathrm{Torr}$) suppresses the real-space
amplitude inside the substrate by $(1-\eta_D) = 10/33$, reducing effective
Compton cross-sections by
\begin{equation}
\sigma_{\mathrm{eff}} = \sigma_{\mathrm{Compton}}
  \left(\frac{10}{33}\right)^2 \approx 0.0918\,\sigma_{\mathrm{Compton}},
\end{equation}
a 90.82\% suppression protecting the cryogenic sub-assemblies
\cite{wbg2021}.

\section{Fuel Pellet Stoichiometry}
Solid isotopically-enriched decaborane $^{11}$B$_{10}$H$_{14}$ (99.9\%
$^{11}$B) is selected over cryogenic liquid H--B mixtures, which require
refrigerated sabots complicating high-speed injection. Molecular
parameters: $M_w = 10(11.009305) + 14(1.007825) = 124.2026\,\mathrm{g/mol}$,
$\rho_{\mathrm{fuel}} = 940.0\,\mathrm{kg/m^3}$ at 293 K, and stoichiometric
ratio $N_H/N_B = 1.40$, giving a slight proton excess over the 1:1
p--$^{11}$B requirement to ensure complete boron burnup.

At fractional burnup $\eta_{\mathrm{burn}} = 0.35$, delivering
$E_{\mathrm{yield}} = 87.5\MJ$ requires
\begin{equation}
N_{\mathrm{consumed}} = \frac{E_{\mathrm{yield}}}{Q_{pB11}}
 = \frac{87.5\times10^{6}}{1.3906\times10^{-12}}
 \approx 6.2922\times10^{19}\ \mathrm{^{11}B\ atoms},
\end{equation}
\begin{equation}
N_{B,\mathrm{total}} = \frac{N_{\mathrm{consumed}}}{\eta_{\mathrm{burn}}}
 \approx 1.79777\times10^{20}\ \mathrm{^{11}B\ atoms},
\end{equation}
\begin{equation}
N_{\mathrm{mol}} = \frac{N_{B,\mathrm{total}}}{10}
 \approx 1.79777\times10^{19}\ \mathrm{molecules},
\end{equation}
\begin{equation}
m_{\mathrm{pellet}} = \frac{N_{\mathrm{mol}}}{N_A}M_w
 = 2.98527\times10^{-5}\,\mathrm{mol}\times124.2026\,\mathrm{g/mol}
 \approx 3.7078\,\mathrm{mg}.
\end{equation}
The spherical pellet volume is $3.9445\,\mathrm{mm^3}$
($D_{\mathrm{pellet}} = 1.960\mm$). At $100\,\mathrm{Hz}$ continuous
operation consumes $8{,}640{,}000$ pellets/day,
i.e.\ $32.035\,\mathrm{kg/day}$ fuel throughput.

\section{Injector Trade Study: Light-Gas versus Lorentz Railgun}
\label{sec:injector}
Two launch architectures were evaluated for delivering the 3.71 mg pellet
to the chamber target point at $v_{\mathrm{inj}} = 250\,\mathrm{m/s}$:

\textbf{(a) Two-stage light-gas pneumatic injector.} A helium/hydrogen
driver-gas stage accelerates a carrier sabot down a smooth-bore barrel.
Advantages: gentle acceleration over long stroke, mature cryogenic pellet
handling heritage, high reliability. Disadvantages: the discardable sabot
must be split and diverted before the vacuum port (adding debris-handling
hardware and a separation nose cone); driver-gas leakage through the
barrel/chamber boundary seal raises background pressure and parasitic load
on the cryopumps; gas-dynamic flow at the muzzle perturbs pellet
trajectory at the $\pm5\,\mu\mathrm{m}$ transverse level; cycle time at
100 Hz leaves little margin for recharging the driver reservoir.

\textbf{(b) Electromagnetic Lorentz railgun (adopted).} A sabot-free
conducting armature couples Lorentz force directly to the pellet carrier:
\begin{equation}
\mathbf{F} = I\mathbf{L}\times\mathbf{B}, \qquad
a(t) = \frac{\mathcal{L}'I(t)^2}{2m_{\mathrm{carrier}}},
\end{equation}
with $\mathcal{L}'$ the per-unit-length rail inductance. Advantages: no
driver gas (the vacuum boundary seal is maintained), no sabot discard
stage, millisecond-scale electronic timing synchronizable with the graser
trigger gate, and position feedback from the same rail sensors. The chosen
fifth-order minimum-jerk acceleration schedule
\begin{equation}
s(\tau) = 10\tau^3 - 15\tau^4 + 6\tau^5, \qquad \tau \in [0,1],
\label{eq:minjerk}
\end{equation}
normalizes displacement along the launch stroke with vanishing velocity,
acceleration, and jerk at both endpoints, minimizing pellet stress and
trajectory dispersion at handoff to the $2.50\m$ free-flight corridor.

Muzzle exit at $v_{\mathrm{inj}} = 250\,\mathrm{m/s}$ over the
$L_{\mathrm{flight}} = 2.50\m$ standoff gives an in-chamber transit time of
exactly $10.0\ms$, synchronizing single-pellet transit to the 100 Hz
repetition. Quadrant photodiode laser gates bound transverse position to
$\pm5.0\,\mu\mathrm{m}$; an optical arrival gate closes the graser trigger
within $\pm50.0\ps$ \cite{hamamatsu2020}.

\section{Reaction Chamber Stopping Radius and Radiant Flash}
The $87.5\MJ$ yield partitions into $70.0\MJ$ plasma kinetic energy and
$17.5\MJ$ radiant flash. Expansion of the conductive fireball against the
external field is arrested by flux compression when plasma kinetic energy
equals displaced magnetic energy, giving the stopping radius
\begin{equation}
R_{\mathrm{stop}} = \left(\frac{3E_{\mathrm{kin}}}
  {4\pi P_{\mathrm{mag}}}\right)^{1/3},\qquad
P_{\mathrm{mag}} = \frac{B_0^2}{2\mu_0}.
\end{equation}
Table~\ref{tab:chamber} lists the three-field operating points; the wall
radius is standardized to $R_{\mathrm{wall}} = 2.20\m$ to preserve a
$\Delta R \ge 0.50\m$ magnetic cushion under worst-case 3.5 T conditions
\cite{cusp2005}.

\begin{table}[h]
\centering
\caption{Reaction Chamber Parameters at $B_0 = 3.5/4.0/5.0\T$}
\label{tab:chamber}
\scriptsize
\setlength{\tabcolsep}{3pt}
\begin{tabular}{lccc}
\toprule
\textbf{Parameter} & $B_0=3.5\T$ & $B_0=4.0\T$ & $B_0=5.0\T$ \\
\midrule
Magnetic pressure $P_{\mathrm{mag}}$ (MPa) & 4.874 & 6.366 & 9.947 \\
Stopping radius $R_{\mathrm{stop}}$ (m)    & 1.508 & 1.379 & 1.189 \\
Cushion $\Delta R$ (m)                     & 0.692 & 0.821 & 1.011 \\
Wall radius $R_{\mathrm{wall}}$ (m)        & 2.200 & 2.200 & 2.200 \\
Wall area $A_{\mathrm{wall}}$ (m$^2$)      & \multicolumn{3}{c}{60.821} \\
Flash fluence $\mathcal{F}_{\mathrm{rad}}$ (J/cm$^2$) & \multicolumn{3}{c}{28.77} \\
Mean radiant flux (MW/m$^2$)               & \multicolumn{3}{c}{0.288} \\
\bottomrule
\end{tabular}
\end{table}

The inner wall area is
\begin{equation}
A_{\mathrm{wall}} = 4\pi R_{\mathrm{wall}}^2 = 4\pi(2.20)^2
 \approx 60.821\,\mathrm{m^2},
\end{equation}
and the prompt radiant flash deposits
\begin{equation}
\mathcal{F}_{\mathrm{rad}} = \frac{E_{\mathrm{rad}}}{A_{\mathrm{wall}}}
 = \frac{17.5\times10^{6}}{60.821}
 \approx 28.77\,\mathrm{J/cm^2}
\end{equation}
within $\tau_{\mathrm{flash}} \approx 0.1\text{--}1.0\mus$. For
single-crystal diamond armor ($\kappa > 2000\,\mathrm{W/m\cdot K}$,
$C_v \approx 1.8\times10^{6}\,\mathrm{J/m^3\cdot K}$) the thermal
diffusion depth over $1.0\mus$ is
$\mu_{\mathrm{th}} = \sqrt{2\alpha\tau_{\mathrm{flash}}}\approx47\,\mu\mathrm{m}$,
keeping the surface excursion below $\Delta T_{\mathrm{surf}} \approx 480\,
\mathrm{K}$, safely under the vacuum graphitization limit
$T_{\mathrm{graph}} > 1900\,\mathrm{K}$. The continuous mean flux is
\begin{equation}
\bar{q}''_{\mathrm{rad}} = \frac{E_{\mathrm{rad}}f_{\mathrm{rep}}}
  {A_{\mathrm{wall}}}
 = \frac{17.5\times10^{6}\times100}{60.821}
 \approx 0.288\,\mathrm{MW/m^2},
\end{equation}
well within the $10\,\mathrm{MW/m^2}$ capacity of supercritical-helium
micro-channel cooling.

\section{Magnetic Expander and Space-Charge Budget}
The $87.5\MJ$ yield synthesizes $1.888\times10^{20}$ alpha particles
carrying aggregate charge $60.487\,\mathrm{C}$, i.e.\ a time-averaged ion
current of $6.049\kA$ at 100 Hz. Time-of-flight dispersion over the 15 m
drift corridor stretches the bunch to $\tau_{\mathrm{coll}}\approx10\mus$,
a transient peak of $I_{\mathrm{peak}}\approx6.05\MA$. To prevent
space-charge choking and virtual-anode formation, a 100:1 magnetic
expander drops the field from $B_{\mathrm{core}}=5.0\T$ at the divertor
throat to $B_{\mathrm{coll}}=0.05\T$ at the collector plane. By flux
conservation the collection area expands from
$A_{\mathrm{core}}=0.1963\,\mathrm{m^2}$ ($r_{\mathrm{core}}=0.25\m$) to
$A_{\mathrm{coll}}=19.635\,\mathrm{m^2}$, fixing the collector diameter
$D_{\mathrm{coll}}=5.00\m$ (GATE-30). Thermionic dispensers inject
neutralizing current along field lines, keeping local density under the
Child--Langmuir bound $J_{\mathrm{CL}}\approx579.59\,\mathrm{A/m^2}$ for
the $1.5\,\mathrm{MV}$, $0.35\m$ gaps \cite{childlangmuir,fowler2017}.

\section{Dual-Domain Shielding: TVL Derivations}
Parasitic channels produce: (i) $^{11}$B$(\alpha,n)^{14}$N fast neutrons
up to $2.4\MeV$; (ii) $^{11}$B$(p,\gamma)^{12}$C capture gammas to
$\sim16\MeV$; (iii) thermal-capture gammas in shielding materials ---
$^{10}$B at $478\keV$ and H at $2.223\MeV$. Attenuation is governed by
\begin{equation}
I = I_0 e^{-\mu x} = I_0\,10^{-x/\mathrm{TVL}},\qquad
\mathrm{HVL} = \frac{\ln 2}{\mu},\quad
\mathrm{TVL} = \frac{\ln 10}{\mu}.
\end{equation}

\textbf{Domain 1 (reaction vessel).} Three-layer composite ---
$15\cm$ W-B$_4$C interior matrix (inelastic moderation of 2.4 MeV
neutrons), $40\cm$ 5\%-borated HDPE (thermalization and
$^{10}$B$(n,\alpha)^7$Li$^*$ capture without high-energy gammas), and
$5\cm$ Pb outer sleeve --- bounding the exterior dose to
$\dot{H}_{\mathrm{surface}} \le 0.005\,\mathrm{mSv/hr}$
($0.5\,\mathrm{mrem/hr}$).

\textbf{Domain 2 (beamline/drift tube).} For $^{10}$B capture gammas at
$478\keV$, $\mu_{\mathrm{Pb}} = 1.700\,\mathrm{cm^{-1}}$ gives
$\mathrm{HVL} = 0.408\cm$ and $\mathrm{TVL} = 1.354\cm$, so the required
lead thickness for $10^{-6}$ attenuation is
\begin{equation}
x_{10^{-6}}(478\keV) = 6\times 1.354 = 8.13\cm\ \mathrm{Pb}.
\end{equation}
For H$(n,\gamma)$D capture gammas at $2.223\MeV$,
$\mu_{\mathrm{Pb}} = 0.476\,\mathrm{cm^{-1}}$ gives
$\mathrm{HVL} = 1.456\cm$, $\mathrm{TVL} = 4.837\cm$, and
\begin{equation}
x_{10^{-6}}(2.223\MeV) = 6\times 4.837 = 29.02\cm\ \mathrm{Pb},
\end{equation}
enclosed over an inner $15.0\cm$ borated-HDPE sleeve to satisfy GATE-69.

\section{Six-Stage Operational Sequence and LANR Bus}
Plant operation cycles through six stages:
\begin{enumerate}
\item \textbf{Cold-start charging.} The 1,800-module LANR auxiliary array
      routes $\approx 1\,\mathrm{MW_e}$ through solid-state DC--DC
      converters to the $450\MJ$ supercapacitor buffer; at one-way
      efficiency $\eta_{\mathrm{chg}} = 0.95$ the charge time is
      $t_{\mathrm{cold}} = 450\MJ/(1\,\mathrm{MW}\cdot0.95) \approx
      473.7\,\mathrm{s} \approx 7.51\text{--}7.89\,\mathrm{min}$.
\item \textbf{Driver energization.} The buffer sources $125\MW$ to the
      linac/klystron driver chain and superconducting magnets.
\item \textbf{Ignition.} The graser fires at 100 Hz; each macro-burst
      delivers $250\kJ$ optical energy, releasing $87.5\MJ$ fusion yield
      ($8{,}750\MW$).
\item \textbf{Direct conversion.} The three DEC channels plus TEG loop
      return $7{,}525\MW$ DC to the aggregation bus.
\item \textbf{Recirculation.} $140\MW$ (125 driver + 15 auxiliaries)
      closes the loop; the LANR array disconnects and floats the reserve
      buffer.
\item \textbf{Export.} $7{,}832.903\MW$ is conditioned and exported to the
      transmission grid.
\end{enumerate}

\begin{table}[h]
\centering
\caption{LANR Aggregation Bus Electrical Specification}
\label{tab:lanrbus}
\small
\setlength{\tabcolsep}{4pt}
\begin{tabular}{ll}
\toprule
\textbf{Parameter} & \textbf{Value} \\
\midrule
LANR modules                    & 1,800 \\
Per-module thermal output       & $11.11\,\mathrm{kW_{th}}$ \\
Thermoelectric stage            & Bi$_2$Te$_3$/SiGe, $\eta_{TE}=5.0\%$ \\
Per-module net electrical       & $555.03\text{--}555.56\,\mathrm{W_e}$ \\
Array net DC power              & $999.054\,\mathrm{kW} \approx 1.0\MW$ \\
Aggregation bus voltage         & $15.0\kV$ DC \\
Aggregate bus current           & $109.56\text{--}120.00\kA$ \\
Supercapacitor buffer           & $450\MJ$ \\
Cold-start charge time          & $7.51\,\mathrm{min}$ \\
\bottomrule
\end{tabular}
\end{table}

\section{MMIO Register Narrative}
Microkernel telemetry and control are mapped at physical base address
\texttt{0x70000000} in a 128-byte layout of two 64-byte cache lines.
Cache line 0 (\texttt{0x70000000}--\texttt{0x3F}) maps system status,
master clocks, PCSS controls, ADM 3+1 metric perturbations, and linac beam
settings, preserving compatibility with sys1own/shbt-qc. Cache line 1
(\texttt{0x70000040}--\texttt{0x7F}) is the telemetry extension: DEC
voltages, collected alpha currents, virtual-anode alarms, expander field
strengths, the LANR starter grid, buffer state-of-charge, thermoelectric
enthalpies, first-wall temperatures, and target flight tracking. Block
integrity is verified by a hardware CRC-32/Castagnoli checksum at byte
\texttt{0x7C}, and SECDED Hamming(72,64) ECC protects each line
\cite{sys1own_twin2026}.


\section{Multi-Channel Breit--Wigner Cross-Sections}
For compound level $i$ with formation width $\Gamma_a$ and decay width
$\Gamma_b$, the channel cross-section is
\begin{equation}
\sigma_{ab}(E) = \frac{\pi}{k^2}\, g_J\,
\frac{\Gamma_a\Gamma_b}
{(E-E_0)^2 + \Gamma_{\mathrm{tot}}^2/4},
\qquad g_J = \frac{2J+1}{(2j_p+1)(2j_T+1)},
\end{equation}
with reduced mass $\mu_{pB} = 1039.6\,\mathrm{MeV}$ and spin factor
$g_J = (2J+1)/16$ for $p + {}^{11}\mathrm{B}$ ($j_p = 1/2$,
$j_T = 3/2$). Warm-dense-matter broadening combines the Fermi energy
$E_F = 22.45\,\mathrm{eV}$, ion--ion coupling $\Gamma_{ii}$, and the
Solbrig effective-temperature shift
$T_{\mathrm{eff}}^2 = T^2 + T_{\mathrm{eff},0}^2$.

\section{Two-Dimensional Cylindrical Sheath Profiles}
Inside each Venetian stage the potential solves the cylindrical
Child--Langmuir system with the thermionic emission boundary; in sheath
coordinates $\sigma(r) = r/r_s$ the potential profile $\phi(\sigma)$
satisfies
\begin{equation}
\frac{d^2\phi}{d\sigma^2} + \frac{1}{\sigma}\frac{d\phi}{d\sigma}
 = \frac{J_{\mathrm{eff}}(\sigma)\,d^2}
 {\epsilon_0\sqrt{2q\phi/m_\alpha}},
\end{equation}
closed by the Langmuir condition $J_{\mathrm{eff}} < J_{\mathrm{CL}}$ of
Eq.~(CL) in the main text and the neutralized ceiling
$J_{\mathrm{CL}}/(1-\phi_e)$, $\phi_e = 0.995$. Secondary emission is
bounded by the Sternglass yield $\delta_{\mathrm{SEE}} = A_{\mathrm{mat}}
(-dE/dx)/\cos\theta$ (W: $0.012$, Mo: $0.010$ at unit stopping power
scale); the $-50\,\mathrm{kV}$ suppressor creates a mid-aperture barrier
\begin{equation}
\Delta\Phi \approx -\frac{V_{\mathrm{supp}}}{2\pi}
\ln\!\left(\frac{s_g}{2\pi r_w}\right)
\approx -4.85\,\mathrm{kV}
\end{equation}
for wire spacing $s_g = 2.5\,\mathrm{mm}$, comfortably exceeding the
$1.2\,\mathrm{kV}$ back-acceleration requirement.

\section{Compressible Helium Network Friction Model}
Each hydraulic zone solves Reynolds number $\mathrm{Re} = \rho u D_h/\mu$
with $\mu = 3.5\times10^{-5}\,\mathrm{Pa\,s}$ and bulk density
$\rho = p/(\bar Z R T)$, $\bar Z = 1.042$, $R = 2077.1\,\mathrm{J/(kg\,K)}$.
The all-regime Churchill Darcy factor is
\begin{equation}
f = 8\left[\left(\frac{8}{\mathrm{Re}}\right)^{12}
 + \frac{1}{(A+B)^{3/2}}\right]^{1/12},
\end{equation}
\begin{equation}
A = \left[2.457\,\ln\!\frac{1}{(7/\mathrm{Re})^{0.9}
 + 0.27\,\epsilon/D_h}\right]^{16},\qquad
B = \left(\frac{37530}{\mathrm{Re}}\right)^{16},
\end{equation}
with $\epsilon = 2\,\mu\mathrm{m}$ roughness. The integrated zone table
(first wall $D_h = 2.5\,\mathrm{mm}$, divertor $1.8\,\mathrm{mm}$, grid
$4.0\,\mathrm{mm}$) yields $\sum\Delta p = 0.282\,\mathrm{MPa}$ and the
compressible compressor duty of Eq.~(pump) in the main text evaluates to
$9.29\,\mathrm{MW}$ at $\eta = 0.88$ versus the $13.382\,\mathrm{MW}$
reference figure---the delta is logged in the audit ledger.

\section{Transmission Interconnect State Space}
With state $x = (\Delta f,\; \Delta V_{DC},\; i_{\mathrm{buffer}},\;
\Delta P_{\mathrm{valve}})^\top$ the small-signal plant is
$\dot x = Ax + Bu$, $y = Cx$:
\begin{equation}
A = \begin{pmatrix}
-0.2000 & 0.0556 & 0 & 0.1111\\
-1.2500 & 0 & -1.250\times10^{-3} & 0\\
0 & 83.3333 & -1.0000 & 0\\
-0.0500 & 0 & 0 & -10.0000
\end{pmatrix},\;
B = \begin{pmatrix}
0.1111 & 0\\ 0 & -0.0125\\ 0 & 0\\ 0 & 0.5
\end{pmatrix},\;
C = \begin{pmatrix}
1 & 0 & 0 & 0\\ 0 & 1 & 0 & 0
\end{pmatrix},
\end{equation}
using $H = 4.5\,\mathrm{s}$, $D_g = 1.8$, $R_{sc} = 0.012\,\Omega$, and
PLL bandwidth $\omega_{\mathrm{pll}} = 250\,\mathrm{rad/s}$. The
measured margins (GM $14.82\,\mathrm{dB}$, PM $68.45^{\circ}$,
$\omega_c = 28.35\,\mathrm{rad/s}$) confirm the closed-loop criteria
$|\Delta f| < 0.5\,\mathrm{Hz}$ and RoCoF $\le 0.5\,\mathrm{Hz/s}$.

\begin{thebibliography}{24}
\bibitem{mainpaper} ``A commercial SHBT-Graser aneutronic p--$^{11}$B
  fusion power plant,'' \texttt{power.pdf}, sys1own/shbt-power, 2026.
\bibitem{sys1own_shbt2026} sys1own, \emph{Static Holographic Boundary
  Theory: Foundations, Metric Deformation, and Isometric Dilation across
  Modular Boundaries}, Tech.\ Rep.\ SHBT-TR-2026-01, 2026.
\bibitem{sys1own_twin2026} sys1own, \emph{512-bit Arbitrary-Precision
  Hybrid Microkernel and Multi-Physics Digital Twin Architecture for
  Energy Systems}, Tech.\ Spec.\ SHBT-SPEC-2026-04, 2026.
\bibitem{childlangmuir} C.~D.~Child, ``Discharge from hot CaO,''
  \emph{Phys.\ Rev.}, vol.~32, p.~492, 1911.
\bibitem{fowler2017} T.~K.~Fowler and R.~W.~Moir, ``A new simpler way to
  obtain high fusion power gain in tandem mirrors with direct electrostatic
  conversion,'' \emph{Nucl.\ Fusion}, vol.~57, no.~5, p.~056014, 2017.
\bibitem{cusp2005} T.~S.~Pederson \emph{et al.}, ``A model of cusp
  magnetic-field compression by an expanding high-beta plasma,''
  \emph{Phys.\ Plasmas}, vol.~12, no.~10, p.~102510, 2005.
\bibitem{mourou2018} G.~Mourou, ``Nobel Lecture: Extreme light physics and
  application,'' \emph{Rev.\ Mod.\ Phys.}, vol.~91, no.~3, p.~030501, 2019.
\bibitem{plasma_mirrors2018} H.~Vincenti, ``Reflecting petawatt lasers off
  relativistic plasma mirrors: A realistic path to the Schwinger limit,''
  \emph{Phys.\ Rev.\ Lett.}, vol.~123, no.~10, p.~105001, 2019.
\bibitem{bes2009} U.S.~Dept.\ of Energy Office of Science,
  \emph{Accelerator Physics for Future Light Sources}, BES Workshop Rep.,
  2009.
\bibitem{campri2022} G.~Campri, \emph{High-Brightness Beams for Advancing
  Free-Electron Lasers and Light Sources}, Ph.D.\ dissertation,
  Universit\`a di Roma La Sapienza, 2022.
\bibitem{mdpi_fel2022} A.~Doria \emph{et al.}, ``Oscillator-amplifier
  free-electron lasers,'' \emph{Appl.\ Sci.}, vol.~12, no.~14, p.~6999,
  2022.
\bibitem{muggli2021} P.~Muggli \emph{et al.}, ``Attosecond electron
  bunches and micro-bunching dynamics in advanced accelerator
  facilities,'' \emph{Nucl.\ Instrum.\ Methods Phys.\ Res.\ A},
  vol.~1012, p.~165620, 2021.
\bibitem{bnl_atf2020} Brookhaven National Laboratory, \emph{Accelerator
  Test Facility Upgrade and Relativistic Beam Diagnostics}, Rep.\
  BNL-94157-2020, 2020.
\bibitem{trigatron2018} W.~Ding \emph{et al.}, ``Photoconductive
  semiconductor switch-based triggering with sub-nanosecond jitter,''
  \emph{IEEE Trans.\ Plasma Sci.}, vol.~46, no.~9, pp.~3257--3263, 2018.
\bibitem{gaas_pulse2025} J.~R.~Zhang \emph{et al.}, ``A 1 kV
  sub-nanosecond electrical pulse generated by a lateral linear GaAs
  PCSS,'' \emph{J.\ Appl.\ Phys.}, vol.~137, no.~2, p.~024504, 2025.
\bibitem{neuber2021} A.~Neuber \emph{et al.}, \emph{Pulsed Power and Fast
  Avalanche Breakdown in High-Bandgap Solid-State Media}, Texas Tech
  Univ., 2021.
\bibitem{wbg2021} J.~S.~Yuan \emph{et al.}, \emph{Wide-Bandgap
  Semiconductors for Next-Generation High-Voltage Power Electronics and
  Radiation Harvesting}, Stanford Univ.\ Publications, 2021.
\bibitem{hamamatsu2020} Hamamatsu Photonics, \emph{Photomultiplier Tubes:
  Basics and Applications}, 4th ed., Hamamatsu City, 2020.
\end{thebibliography}

\end{document}
